\begin{textsample}{POS Dim 3 – mg – Score 18.00 – mg/mg\_inf\_campo\_exp\_6.txt}  \label{ex:f3_pos_242}
2.6.5 Campo de experiências : Espaços, tempos, quantidades, relações e transformações As \textbf{crianças} vivem inseridas em espaços e tempos de diferentes dimensões, em um mundo constituído de fenômenos naturais e socioculturais.
Desde muito \textbf{pequenas}, elas procuram se situar em diversos espaços ( rua, bairro, cidade etc. ) e tempos ( dia e noite ; hoje, ontem e amanhã etc. ). \textbf{Demonstram} também \textbf{curiosidade} sobre o mundo físico ( seu próprio corpo, os fenômenos atmosféricos, os animais, as plantas, as transformações da natureza, os diferentes tipos de materiais e as possibilidades de sua manipulação etc. ) e o mundo sociocultural ( as relações de parentesco e sociais entre as pessoas que conhece ; como vivem e em que trabalham essas pessoas ; quais suas tradições e seus costumes ; a diversidade entre elas etc. ). [... ] as \textbf{crianças} também se deparam, \textbf{frequentemente} com conhecimentos matemáticos ( contagem, ordenação, relações entre quantidades, dimensões, medidas, comparação de \textbf{pesos} e de comprimentos, avaliação de distâncias, reconhecimento de formas geométricas, conhecimento e reconhecimento de numerais cardinais e ordinais etc. ), que igualmente aguçam a \textbf{curiosidade} ( BRASIL, 2017, p.40-41 ).
As instituições de Educação \textbf{Infantil} devem promover experiências para que as \textbf{crianças} construam a percepção de espaços, tempos, quantidades, relações e transformações presentes no seu dia a dia motivando -as a terem um olhar mais crítico e criativo do mundo.
Para que isso ocorra é preciso valorizar todo o conhecimento que a \textbf{criança} traz de suas práticas sociais cotidianas, pois toda \textbf{criança} apresenta uma diversidade de conhecimentos que pode e deve servir de ponto de partida para novas aprendizagens.
Além disso, é preciso que ela atribua sentido às experiências mediadas pelo docente para que elabore hipóteses na resolução de problemas contextualizados, teste a validade dessas hipóteses, e modifique -as para construir e reconstruir novos conhecimentos.
Os \textbf{bebês} e as demais \textbf{crianças} estão inseridos numa cultura que lida constantemente com noções que envolvem esse campo de experiência.
Vivenciam situações em que as pessoas utilizam a contagem para resolver problemas \textbf{diários}, fazem pagamentos e trocos, comparam tamanhos, calculam os dias que faltam para uma determinada data, \textbf{brincam} ao telefone, trocam canais na televisão, exploram o espaço, recitam a sucessão dos números, entre tantas outras atividades corriqueiras do dia a dia.
Imersas em um meio repleto de produtos da cultura, as \textbf{crianças} ao \textbf{manipular} objetos e outros materiais, agem para entender seu funcionamento, para diferenciar suas características, cada vez mais com as frequentes perguntas “ como? ” e “ por quê? ” dirigidas a parceiros mais experientes. ( BRASIL, 2018 ).
Quanto tempo falta para o meu aniversário?
Por que quando minha avó era \textbf{criança} não havia televisão?
Por que alguns objetos afundam e outros não?
Por que existem alguns animais com penas e outros com pelos?
Quantas vezes um elefante é maior do que um cavalo?
Como estes doces podem ser distribuídos igualmente entre as \textbf{crianças}?
Que jogador de futebol fez mais gols na Copa?
Uma centopeia tem mais patas do que uma abelha?
A forma como a instituição de Educação \textbf{Infantil} e o docente tratam essa \textbf{curiosidade}, esses questionamentos, esse \textbf{desejo} de também se apropriar desse saber pode lhes ajudar ( ou não ) a observar regularidades e permanências, ou diversidades e mudanças na natureza e na vida social, a formular noções de espaço, de tempo e a fazer aproximações em torno da ideia de causalidade e transformação. ( BRASIL, 2018 ).
À medida que a \textbf{professora} considera as unidades de Educação \textbf{Infantil} como ambientes onde a \textbf{curiosidade} das \textbf{crianças} sobre o mundo físico e social pode alimentar a construção por elas de noções, comparações e implicações, ela as \textbf{ajuda} a construir explicações, conforme percebe seus \textbf{gestos}, sentimentos, intuições, seus motivos e sentidos pessoais nas respostas que elas dão.
Nesse processo procuram articular o modo como as \textbf{crianças} agem, sentem e pensam com os conhecimentos já disponíveis na cultura sobre cada objeto de conhecimento. ( BRASIL, 2018, p.93 ).
Na interação e observação das \textbf{crianças}, o docente pode auxiliar as \textbf{crianças} a perceberem relações entre objetos e materiais, chamar -lhes a atenção para certos aspectos da situação, estimulá -las a fazer novas descobertas e construir novos conhecimentos a partir dos saberes que já possuem.
O docente deve lidar com as opiniões \textbf{infantis}, acompanhando o constante esforço que cada \textbf{criança} faz para singularizar -se, e convidando -a, num processo contínuo, a construir uma sociedade mais solidária, inclusiva, equitativa e igualitária.
As DCNEI/2009, no seu artigo 9º, destacam a importância de incentivar a \textbf{curiosidade}, a exploração, o encantamento, o questionamento e a indagação das \textbf{crianças} sobre o mundo que as cercam.
O mesmo artigo destaca que devem ser garantidas às \textbf{crianças} experiências que “ promovam a interação, o \textbf{cuidado}, a preservação e o conhecimento da biodiversidade e da sustentabilidade da vida na Terra, assim como o não \textbf{desperdício} dos recursos naturais ” ( BRASIL, 2009 ).
É importante oferecer às \textbf{crianças} desde \textbf{pequenas}, oportunidades de explorar diferentes tipos de objetos, seres e materiais da natureza, fenômenos físicos, químicos e biológicos, bem como o meio ambiente e sua sustentabilidade.
As \textbf{crianças} vão formulando e reformulando hipóteses e ideias explicativas sobre o mundo que as cercam, vão querendo saber “ o porquê ” das coisas, conhecendo -as, transformando -as e fazendo novas invenções.
Portanto, a instituição de Educação \textbf{Infantil} deve promover junto às \textbf{crianças} constantes atividades exploratórias, contato com materiais diversos que permitam as \textbf{crianças}, através das interações e \textbf{brincadeiras}, pensar os mundos da natureza e da sociedade, incluindo os animais, as plantas, os objetos, a tecnologia, o comportamento humano e outros aspectos da cultura, assim, elas vivenciam de modo integrado, experiências em relação ao tempo, ao espaço, às quantidades, relações e transformações.

% matched lemmas: ajuda, bebê, brincadeira, brincar, criança, cuidado, curiosidade, demonstrar, desejo, desperdício, diário, frequentemente, gesto, infantil, manipular, pequeno, peso, professora
\end{textsample}
